\section{Einleitung}
\pagenumbering{arabic}
Das Graph-Burning-Problem ist ein kombinatorisches Optimierungsproblem der theoretischen Informatik. 
Das Modell beschreibt einen zeitdiskreten Prozess, bei dem in jedem Schritt ein neuer Brandherd gelegt wird, während die aktive Feuerfront gleichzeitig auf die Nachbarknoten im Graphen übergreift. Das Ziel besteht darin, alle Knoten des Graphen in einer minimalen Anzahl von Zeitschritten zu verbrennen \cite{NAZERI2023}.

Das Problem ist für viele Anwendungen relevant. Es kann beispielsweise eingesetzt werden, um eine epidemiologische Ausbreitung \cite{Rosh16, GLM21}, einen Virusbefall in Computernetzwerken \cite{Rosh16}, die Verbreitung von Informationen in sozialen Netzwerken \cite{bjr14, Rosh16, Gau22, NAZERI2023, Ashok25} oder die Alarmierung der Bevölkerung \cite{Si19, SLP19, GLM21} im Zeitablauf zu modellieren. 
Ziel dieser Arbeit ist es, den gewählten Ansatz zu implementieren und unter Verwendung von verschiedenen Zentralitätsmaßen und Penalty-Funktionen empirisch zu untersuchen. 

Nach der Einleitung werden in Kapitel~\ref{cap:theo} die theoretischen Grundlagen vorgestellt. 
In Kapitel~\ref{cap:nazeri} wird dann der zu optimierende Ansatz von Nazeri et al. \cite{NAZERI2023} beschrieben. 
In Kapitel~\ref{cap:impl} folgt die Erläuterung der eigenen Implementierung. 
In Kapitel~\ref{cap:num} werden die numerischen Ergebnisse dargestellt. 
Die Arbeit wird in Kapitel~\ref{cap:conc} mit einem Fazit abgeschlossen. 
